% ═══════════════════════════════════════════════════════════════
% BEAMER: DS-01 - Mi nueva casa
% Klasse 9, Unidad 2
% ═══════════════════════════════════════════════════════════════

\documentclass[aspectratio=169, 11pt]{beamer}

\usepackage[utf8]{inputenc}
\usepackage[T1]{fontenc}
\usepackage[ngerman]{babel}

\usepackage{helvet}
\renewcommand{\familydefault}{\sfdefault}
\usepackage{palatino}
\newcommand{\seriftext}[1]{{\fontfamily{ppl}\selectfont #1}}

\usepackage{tikz}
\usepackage{array}
\usepackage{booktabs}
\usepackage{pifont}

% ═══════════════════════════════════════════════════════════════
% FARBEN
% ═══════════════════════════════════════════════════════════════

\definecolor{akzent}{HTML}{C41E3A}
\definecolor{stamm}{HTML}{C62828}
\definecolor{hellgrau}{HTML}{F5F5F5}
\definecolor{dunkelgrau}{HTML}{333333}
\definecolor{mittelgrau}{HTML}{666666}
\definecolor{richtig}{HTML}{2A7A4B}
\definecolor{falsch}{HTML}{C62828}

% ═══════════════════════════════════════════════════════════════
% BEAMER-THEME
% ═══════════════════════════════════════════════════════════════

\usetheme{default}
\usecolortheme{default}

\setbeamertemplate{navigation symbols}{}

\setbeamercolor{frametitle}{fg=white, bg=akzent}
\setbeamercolor{title}{fg=white}
\setbeamercolor{subtitle}{fg=white}
\setbeamercolor{author}{fg=white}
\setbeamercolor{date}{fg=white}
\setbeamercolor{structure}{fg=akzent}
\setbeamercolor{itemize item}{fg=akzent}
\setbeamercolor{itemize subitem}{fg=mittelgrau}
\setbeamercolor{block title}{fg=white, bg=akzent}
\setbeamercolor{block body}{bg=hellgrau}

\setbeamertemplate{itemize item}{\textbullet}
\setbeamertemplate{itemize subitem}{--}

\setbeamertemplate{title page}{
    \begin{tikzpicture}[remember picture, overlay]
        \fill[akzent] (current page.north west) rectangle (current page.south east);
    \end{tikzpicture}
    \vspace{2cm}
    \begin{center}
        {\Huge\bfseries\textcolor{white}{\inserttitle}}\\[0.8cm]
        {\Large\textcolor{white}{\insertsubtitle}}\\[1.5cm]
        {\large\textcolor{white}{\insertauthor}}\\[0.3cm]
        {\textcolor{white}{\insertdate}}
    \end{center}
}

\setbeamertemplate{frametitle}{
    \nointerlineskip
    \begin{beamercolorbox}[wd=\paperwidth, ht=1.2cm, dp=0.3cm]{frametitle}
        \hspace{0.5cm}\Large\bfseries\insertframetitle
    \end{beamercolorbox}
}

\setbeamertemplate{footline}{
    \hfill\textcolor{mittelgrau}{\footnotesize\insertframenumber}\hspace{0.5cm}\vspace{0.3cm}
}

% ═══════════════════════════════════════════════════════════════
% BEFEHLE
% ═══════════════════════════════════════════════════════════════

\newcommand{\es}[1]{\seriftext{\textbf{#1}}}
\newcommand{\de}[1]{\textit{\textcolor{mittelgrau}{#1}}}
\newcommand{\stw}[1]{\textcolor{stamm}{\textbf{#1}}}
\newcommand{\abmark}[1]{{\Large$\rightarrow$ \textbf{Arbeitsblatt Aufgabe #1}}}
\newcommand{\richmark}{\textcolor{richtig}{\ding{51}}}
\newcommand{\falschmark}{\textcolor{falsch}{\ding{55}}}

% ═══════════════════════════════════════════════════════════════
% METADATEN
% ═══════════════════════════════════════════════════════════════

\title{Mi nueva casa}
\subtitle{Das Haus und der Umzug}
\author{Spanisch · Klasse 9}
\date{}

% ═══════════════════════════════════════════════════════════════
% DOKUMENT
% ═══════════════════════════════════════════════════════════════

\begin{document}

% ---------------------------------------------------------------
% TITELFOLIE
% ---------------------------------------------------------------
\begin{frame}[plain]
    \titlepage
\end{frame}

% ---------------------------------------------------------------
% LERNZIELE
% ---------------------------------------------------------------
\begin{frame}{Lernziele}
    \vspace{0.5cm}

    {\large Heute lernst du:}

    \vspace{0.6cm}

    \begin{itemize}
        \setlength{\itemsep}{0.5cm}
        \item \textbf{Hausräume und Möbel benennen}\\
              {\small 7 Räume + 17 Gegenstände}
        \item \textbf{Bewegungsverben mit Präpositionen verwenden}\\
              {\small bajar de, subir a, poner en, sacar de, traer de, llevar a}
        \item \textbf{Objektpronomen richtig einsetzen}\\
              {\small lo, la, los, las}
        \item \textbf{Einen Umzug beschreiben}\\
              {\small Fragen stellen und beantworten}
    \end{itemize}
\end{frame}

% ═══════════════════════════════════════════════════════════════
% BLOCK 1: HAUSRÄUME
% ═══════════════════════════════════════════════════════════════

\begin{frame}{Las partes de la casa}
    \vspace{0.3cm}

    {\large Welche Räume hat eine Wohnung?}

    \vspace{0.5cm}

    \begin{columns}[T]
        \begin{column}{0.45\textwidth}
            \begin{tabular}{@{}ll@{}}
                \es{el pasillo} & \de{der Flur} \\[0.4em]
                \es{el dormitorio} & \de{das Schlafzimmer} \\[0.4em]
                \es{el salón} & \de{das Wohnzimmer} \\[0.4em]
                \es{la cocina} & \de{die Küche} \\
            \end{tabular}
        \end{column}

        \begin{column}{0.45\textwidth}
            \begin{tabular}{@{}ll@{}}
                \es{el balcón} & \de{der Balkon} \\[0.4em]
                \es{el baño} & \de{das Badezimmer} \\[0.4em]
                \es{el aseo} & \de{die Toilette} \\
            \end{tabular}
        \end{column}
    \end{columns}

    \vspace{0.8cm}

    {\small\textit{Tipp:} \es{baño} = Bad mit Dusche, \es{aseo} = nur WC}
\end{frame}

% ---------------------------------------------------------------
\begin{frame}{Los muebles y objetos}
    \vspace{0.2cm}

    {\large Was steht in den Zimmern?}

    \vspace{0.3cm}

    \begin{columns}[T]
        \begin{column}{0.3\textwidth}
            {\small
            \begin{tabular}{@{}ll@{}}
                \es{la lámpara} & Lampe \\[0.2em]
                \es{la cama} & Bett \\[0.2em]
                \es{el armario} & Schrank \\[0.2em]
                \es{la estantería} & Regal \\[0.2em]
                \es{la planta} & Pflanze \\[0.2em]
                \es{el sillón} & Sessel \\
            \end{tabular}}
        \end{column}

        \begin{column}{0.3\textwidth}
            {\small
            \begin{tabular}{@{}ll@{}}
                \es{el sofá} & Sofa \\[0.2em]
                \es{el cuadro} & Bild \\[0.2em]
                \es{el televisor} & Fernseher \\[0.2em]
                \es{la nevera} & Kühlschrank \\[0.2em]
                \es{el cubo de basura} & Mülleimer \\[0.2em]
                \es{el reloj} & Uhr \\
            \end{tabular}}
        \end{column}

        \begin{column}{0.3\textwidth}
            {\small
            \begin{tabular}{@{}ll@{}}
                \es{el horno} & Ofen \\[0.2em]
                \es{la ducha} & Dusche \\[0.2em]
                \es{el espejo} & Spiegel \\[0.2em]
                \es{la guitarra} & Gitarre \\[0.2em]
                \es{el póster} & Poster \\
            \end{tabular}}
        \end{column}
    \end{columns}
\end{frame}

% --- APPLY ---
\begin{frame}{Jetzt du!}
    \vspace{1cm}

    \begin{center}
        \abmark{1}

        \vspace{1cm}

        {\large Ordne die Bilder den spanischen Wörtern zu.}

        \vspace{0.5cm}

        {\small Zeit: 10 Minuten}
    \end{center}
\end{frame}

% ═══════════════════════════════════════════════════════════════
% BLOCK 2: BEWEGUNGSVERBEN
% ═══════════════════════════════════════════════════════════════

\begin{frame}{Los verbos de movimiento}
    \vspace{0.3cm}

    {\large Bei einem Umzug brauchst du diese Verben:}

    \vspace{0.5cm}

    \begin{center}
    \begin{tabular}{@{}ccl@{}}
        $\downarrow$ & \es{bajar de} & \de{herunterbringen von} \\[0.5em]
        $\uparrow$ & \es{subir a} & \de{hochbringen nach} \\[0.5em]
        $\oplus$ & \es{poner en} & \de{stellen/legen in} \\[0.5em]
        $\ominus$ & \es{sacar de} & \de{herausnehmen aus} \\[0.5em]
        $\leftarrow$ & \es{traer de} & \de{bringen (her) von} \\[0.5em]
        $\rightarrow$ & \es{llevar a} & \de{bringen (hin) nach} \\
    \end{tabular}
    \end{center}
\end{frame}

% ---------------------------------------------------------------
\begin{frame}{Die Präpositionen}
    \vspace{0.5cm}

    \begin{block}{Merke: Richtung}
        \textbf{de / del} = von, aus (Richtung \textbf{weg})\\
        \textbf{a / al} = nach, zu (Richtung \textbf{hin})\\
        \textbf{en} = in (Position, \textbf{wo?})
    \end{block}

    \vspace{0.5cm}

    \textbf{Beispiele:}

    \vspace{0.3cm}

    \es{Bajo la caja \textbf{del} camión.}\\
    \de{Ich bringe die Kiste aus dem LKW.}

    \vspace{0.3cm}

    \es{Subo los muebles \textbf{al} piso.}\\
    \de{Ich bringe die Möbel in die Wohnung hoch.}
\end{frame}

% --- APPLY ---
\begin{frame}{Jetzt du!}
    \vspace{1cm}

    \begin{center}
        \abmark{2}

        \vspace{1cm}

        {\large Ergänze die Verben und Präpositionen.}

        \vspace{0.5cm}

        {\small Zeit: 8 Minuten}
    \end{center}
\end{frame}

% ═══════════════════════════════════════════════════════════════
% BLOCK 3: OBJEKTPRONOMEN
% ═══════════════════════════════════════════════════════════════

\begin{frame}{Los pronombres de objeto directo}
    \vspace{0.3cm}

    {\large Statt das Nomen zu wiederholen, nutzt du Pronomen:}

    \vspace{0.5cm}

    \begin{center}
    \begin{tabular}{@{}lcc@{}}
        & \textbf{Singular} & \textbf{Plural} \\[0.3em]
        \midrule
        maskulin & \es{lo} \de{(ihn)} & \es{los} \de{(sie)} \\[0.4em]
        feminin & \es{la} \de{(sie)} & \es{las} \de{(sie)} \\
    \end{tabular}
    \end{center}

    \vspace{0.5cm}

    \begin{block}{Stellung}
        Das Pronomen steht \textbf{vor} dem konjugierten Verb!
    \end{block}
\end{frame}

% ---------------------------------------------------------------
\begin{frame}{Beispiele}
    \vspace{0.5cm}

    {\large Nomen $\rightarrow$ Pronomen}

    \vspace{0.8cm}

    \es{¿La mesa?} $\rightarrow$ \es{La ponemos en el salón.}\\
    \de{Der Tisch? Wir stellen ihn ins Wohnzimmer.}

    \vspace{0.6cm}

    \es{¿Los muebles?} $\rightarrow$ \es{Los subimos al piso.}\\
    \de{Die Möbel? Wir bringen sie in die Wohnung hoch.}

    \vspace{0.6cm}

    \es{¿Las plantas?} $\rightarrow$ \es{Las traemos del camión.}\\
    \de{Die Pflanzen? Wir bringen sie aus dem LKW.}
\end{frame}

% --- APPLY ---
\begin{frame}{Jetzt du!}
    \vspace{1cm}

    \begin{center}
        \abmark{3}

        \vspace{1cm}

        {\large Ersetze das Nomen durch ein Pronomen.}

        \vspace{0.5cm}

        {\small Zeit: 5 Minuten}
    \end{center}
\end{frame}

% ═══════════════════════════════════════════════════════════════
% BLOCK 4: REDEMITTEL + DIALOG
% ═══════════════════════════════════════════════════════════════

\begin{frame}{Comunicación: La mudanza}
    \vspace{0.3cm}

    {\large Wie fragst du beim Umzug?}

    \vspace{0.5cm}

    \begin{columns}[T]
        \begin{column}{0.45\textwidth}
            \textbf{Fragen:}

            \vspace{0.3cm}

            \es{¿Dónde ponemos el/la...?}\\
            \de{Wo stellen wir den/die... hin?}

            \vspace{0.3cm}

            \es{¿Dónde está el/la...?}\\
            \de{Wo ist der/die...?}
        \end{column}

        \begin{column}{0.45\textwidth}
            \textbf{Antworten:}

            \vspace{0.3cm}

            \es{Lo/La ponemos en...}\\
            \de{Wir stellen ihn/sie in...}

            \vspace{0.3cm}

            \es{Está al lado de / en...}\\
            \de{Er/Sie ist neben / in...}
        \end{column}
    \end{columns}
\end{frame}

% --- APPLY ---
\begin{frame}{Jetzt du!}
    \vspace{0.8cm}

    \begin{center}
        \abmark{4}

        \vspace{0.8cm}

        {\large Arbeitet zu zweit:}

        \vspace{0.3cm}

        {\normalsize Fragt euch gegenseitig, wo ihr die Möbel hinstellt.\\
        Verwendet \es{lo/la/los/las}!}

        \vspace{0.5cm}

        {\small Zeit: 8 Minuten}
    \end{center}
\end{frame}

% ═══════════════════════════════════════════════════════════════
% ABSCHLUSS
% ═══════════════════════════════════════════════════════════════

\begin{frame}{Zusammenfassung}
    \vspace{0.3cm}

    \begin{columns}[T]
        \begin{column}{0.45\textwidth}
            \textbf{Vokabeln}

            \begin{itemize}
                \item 7 Hausräume
                \item 17 Möbel/Gegenstände
            \end{itemize}

            \vspace{0.3cm}

            \textbf{Grammatik}

            \begin{itemize}
                \item 6 Bewegungsverben + Präp.
                \item 4 Objektpronomen
            \end{itemize}
        \end{column}

        \begin{column}{0.45\textwidth}
            \textbf{Die wichtigsten Sätze:}

            \vspace{0.3cm}

            \es{¿Dónde ponemos...?}

            \vspace{0.2cm}

            \es{Lo/La ponemos en...}

            \vspace{0.5cm}

            \textbf{Hausaufgabe:}

            \vspace{0.3cm}

            Vokabeln lernen (24 Wörter)
        \end{column}
    \end{columns}
\end{frame}

% ---------------------------------------------------------------
\begin{frame}{Typische Fehler vermeiden}
    \vspace{0.5cm}

    \begin{center}
        \begin{tabular}{@{}cl@{\hspace{1.5cm}}cl@{}}
            \textcolor{falsch}{\Large$\times$} & \es{La ponemos \stw{a} el salón.} &
            \textcolor{richtig}{\Large$\checkmark$} & \es{La ponemos \stw{en} el salón.} \\[0.8em]

            \textcolor{falsch}{\Large$\times$} & \es{Bajo \stw{en} el camión.} &
            \textcolor{richtig}{\Large$\checkmark$} & \es{Bajo \stw{del} camión.} \\[0.8em]

            \textcolor{falsch}{\Large$\times$} & \es{Ponemos \stw{lo} mesa...} &
            \textcolor{richtig}{\Large$\checkmark$} & \es{Ponemos \stw{la} mesa...} \\
        \end{tabular}
    \end{center}

    \vspace{0.5cm}

    {\small\textit{Merke:} \es{poner en}, \es{bajar de}, \es{subir a} -- Präposition beachten!}
\end{frame}

\end{document}
